\documentclass[10pt,a4paper]{amsart}
\usepackage[margin=19mm]{geometry}
\usepackage{amsmath,amssymb,mathtools}
\usepackage[scr=rsfs,cal=euler]{mathalfa}
\usepackage{xfrac}
\usepackage[unicode,hidelinks,bookmarks=false]{hyperref}
\newcommand{\ie}{i.e.\ }
\newcommand{\cf}{cf.\ }
\newcommand{\resp}{respectively\ }
\newcommand{\Subsection}{\subsection}
\providecommand{\nobreakdash}{\nobreak\hbox{-}}
\setlength{\emergencystretch}{2em}
\multlinegap0pt
\newtheorem{theorem}{Theorem}
\newtheorem{conjecture}[theorem]{Conjecture}
\newtheorem{problem}[theorem]{Problem}
\newtheorem{example}[theorem]{Example}
\newcommand{\C}{\mathbb{C}}
\newcommand{\R}{\mathbb{R}}
\newcommand{\T}{\mathbb{T}}
\newcommand{\hatM}{\widehat{M}}
\newcommand{\calL}{\mathcal{L}}
\begin{document}
\title[Holomorphic Hulls for Compact 3-Manifolds]{Holomorphic Hulls for Compact 3-Manifolds}
\author{Ali M. Elgindi}
\date{}
\maketitle
\noindent\emph{Source and licence.} Holomorphic Hulls for Compact 3-Manifolds. Version arXiv:2606.24951v1 (CC BY 4.0). This material is distributed under the Creative Commons Attribution 4.0 International licence, http://creativecommons.org/licenses/by/4.0/, as recorded in the arXiv OAI licence metadata for this exact version and shown on the abstract page https://arxiv.org/abs/2606.24951v1. Full rights notice: licenses/arxiv-2606.24951-CC-BY-4.0.txt.\par\smallskip
\noindent\textit{Editorial scope.} Counted unit: the original \texttt{theorem} environment \texttt{-} at source lines 90--95 together with its complete proof at lines 96--98; the delivered document keeps source lines 77--102 so that every referenced label and every citation resolves inside it. Native editable source is the paper's own concatenated TeX; the AMS wrapper is editorial formatting only. No publisher class, upstream script or unrelated artwork is redistributed.
\section{Main Results}
\subsection{Elliptic Components}
\begin{theorem}[Elliptic Components] Let $M \hookrightarrow \mathbb{C}^3$ be an embedding as in [1] with a link component $K \subset L$ where $\theta > \frac{\pi}{2}$ gives elliptic type with Bishop invariant $\gamma < \frac{1}{2}$. Then the holomorphic hull $\hatM$ contains a smooth Levi-flat hypersurface $\mathcal{H}$ with its boundary lying on the knot of complex tangents: $\partial\mathcal{H} \subset K$ and $K$ is itself foliated by analytic discs whose boundaries lie in $M$.
\end{theorem}
\begin{proof}
By Bishop's theorem in [3], for each elliptic complex tangent point: $x \in K, \gamma (x)< \frac{1}{2}$) there exists a one-parameter family of analytic discs ${D_{x,t},\{t \in (0,\epsilon]}\}$ with boundaries in $M$ converging to $x$. Kenig and Webster [8] proved that the union of these discs near $x$ is a smooth Levi-flat hypersurface. Since the embedding in [4] varies smoothly along $x \in K$, these local families glue together to form a global Levi-flat hypersurface $\mathcal{H}$ with its boundary $\partial \mathcal{H} \subset K$. The discs foliate $\mathcal{H}$ by construction.
\end{proof}
\subsection{Hyperbolic Components}
\begin{theorem}[Hyperbolic Components]
Let $M \hookrightarrow \mathbb{C}^3$ be as above with a complex tangent link component $K$ where the analytic angle $\theta < \frac{\pi}{2}$, making that the points are of  Bishop hyperbolic type, $\gamma > \frac{1}{2}$. (by our work in [4])  Then there exists a neighborhood U of $K$ such that: $\hat{M} \cap U = M \cap U$; i.e., no other local hull formation occurs.\end{theorem}\begin{proof}Bishop [5] proved that for hyperbolic points ($\gamma > \frac{1}{2}$) no analytic discs attach to $M$ in a neighborhood of the point. The hull $\hat{M}$ is defined as the set of points that cannot be separated from $M$ by holomorphic functions. Without attaching discs, the hull cannot extend beyond $M$ locally. Hence, there exists a neighborhood $U$ of $K$ such that $\widehat{M \cup U} = M \cup U$

\end{proof}
\subsection{Parabolic Components and Jöricke's Onion}

\begin{theorem}
[Parabolic Components]Let $M \hookrightarrow \mathbb{C}^3$ be as above with a link component $K$ where $\theta = \frac{\pi}{2}$ , and as such has Bishop parabolic type: $\gamma = \frac{1}{2}$. The
\begin{enumerate}\item If the Jöricke index vanishes, the hull is locally trivial: $\hat{M} = M$ near $K$
\item If the Jöricke index is positive, the hull exhibits nested Levi-flat shells (an ``onion'') centered along $K$, as described in [6,7].
\end{enumerate}
\end{theorem}

\begin{proof}
Jöricke [6,7] proved that at a parabolic point ($\gamma = \frac{1}{2}$) the local hull structure is determined by a discrete index. If the index is zero, the hull coincides with $M$ in a neighborhood of the point. If the index is $k > 0$, the hull contains $k$ nested Levi-flat nested shells (an onion) centered at the point. By the smoothness of the construction in [1], these local structures vary smoothly along $K$, yielding a global onion structure along the entire component
\end{proof}

\subsection{Transitions Along a Single Knot}
\begin{theorem}[Transitions]Let $K$ be a single knot component of the complex tangent link $L\subset M$ such that the holomorphic angle $\theta$ varies along $K$ (see [4]), taking values both $> \frac{\pi}{2}$ and $< \frac{\pi}{2}$. By our work in[4], this corresponds respectively with the hyperbolic and elliptic Bishop invariant construction ($\gamma>\frac{1}{2}, <\frac{1}{2}$). Furthermore, at points where the holomorphic angle is $\theta = \frac{\pi}{2}$ (corresponding to parabolic Bishop $\gamma=\frac{1}{2}$), the hull $\hat{M}$ develops a corner or singularity\end{theorem}
\begin{proof}On the elliptic side of a parabolic point $p \in K$, Theorem 3.1 gives a Levi-flat tube attached to $K$. On the hyperbolic side, Theorem 3.2 gives no attachment. The boundary of the tube must therefore lie exactly at $p$. This abrupt change in the hull's geometry produces a corner or singularity at $p$. The exact nature of the singularity depends on the Jöricke index at $p$; for index $k = 0$ the corner is sharp, while for $k > 0$ the onion shells smooth the transition.\end{proof} 


\end{document}
